\documentclass[11pt,a4paper]{article}
\usepackage[margin=25mm]{geometry}
\usepackage{fontspec}
\setmainfont{TeX Gyre Pagella}
\usepackage{amsmath,amssymb,mathtools}
\usepackage{graphicx}
\usepackage{xcolor}
\usepackage[unicode,colorlinks=true,linkcolor=blue,urlcolor=blue]{hyperref}
\usepackage{bookmark}
\usepackage{enumitem}
\setlist{itemsep=3pt,topsep=5pt}
\setlength{\emergencystretch}{3em}
\newcommand{\tightlist}{\setlength{\itemsep}{0pt}\setlength{\parskip}{0pt}}
\newcommand{\passthrough}[1]{#1}
\newcommand{\pandocbounded}[1]{#1}
\makeatletter
\def\maxwidth{\ifdim\Gin@nat@width>\linewidth\linewidth\else\Gin@nat@width\fi}
\def\maxheight{\ifdim\Gin@nat@height>0.75\textheight0.75\textheight\else\Gin@nat@height\fi}
\makeatother
\setkeys{Gin}{width=\maxwidth,height=\maxheight,keepaspectratio}
\hypersetup{pdfauthor={OpenAI GPT-6.1 Sol, Codex, Ultra effort}}
\begin{document}
\section{Cyclotomic polynomials and their
automorphisms}\label{cyclotomic-polynomials-and-their-automorphisms}

\emph{Written by GPT-6.1 Sol (OpenAI) in Codex, Ultra setting, October
2026. Self-checked by the writing AI. Independent full-lesson AI review
is pending. Public domain (CC0).}

A supporting excerpt: definitions and the complete Theorem 12.1 proof,
not the full parent lesson.

Fix the complex embedding and write \[
\zeta_n=e^{2\pi i/n},\qquad K_n=\mathbf Q(\zeta_n),\qquad
U_n=(\mathbf Z/n\mathbf Z)^\times.
\tag{1}
\] We take \(U_1\) to be the trivial group and \(\varphi(1)=1\). Both
\(K_1\) and \(K_2\) are \(\mathbf Q\).

\subsection{Cyclotomic polynomials and their
automorphisms}\label{cyclotomic-polynomials-and-their-automorphisms-1}

Define \[
\Phi_n(X)=\prod_{\substack{1\leq a\leq n\\(a,n)=1}}(X-\zeta_n^a).
\] Grouping roots by their exact orders gives \[
X^n-1=\prod_{d\mid n}\Phi_d(X).
\tag{2}
\] Inductively, every \(\Phi_n\) is monic and integral. Indeed, divide
\(X^n-1\) by the product of the already integral monic factors for
proper divisors. Monic division has integer quotient and remainder; the
identity over \(\mathbf C\) makes that remainder zero.

\textbf{Monic factor lemma.} If a monic polynomial in \(\mathbf Z[X]\)
is a product of two monic polynomials \(f,g\in\mathbf Q[X]\), then
\(f,g\in\mathbf Z[X]\).

\textbf{Proof.} After clearing denominators and dividing the resulting
coefficients by their greatest common divisor, write \(f_0=af\),
\(g_0=bg\), where \(f_0,g_0\) are primitive integer polynomials and
\(a,b\) are positive integers: they are the leading coefficients because
\(f,g\) are monic. The product of two primitive polynomials is
primitive. Indeed, for each prime \(p\) their reductions are nonzero
polynomials over \(\mathbf F_p\); the product is nonzero because that
polynomial ring is a domain. Thus no prime divides every coefficient of
the product. But \(f_0g_0=abfg\) has coefficient greatest common divisor
\(ab\), since \(fg\) is monic and integral. Consequently \(ab=1\), so
both factors were integral. \(\square\)

\textbf{Theorem 12.1.} The polynomial \(\Phi_n\) is irreducible over
\(\mathbf Q\). Its degree is \(\varphi(n)\), and the map \[
U_n\xrightarrow{\sim}\operatorname{Gal}(K_n/\mathbf Q),
\qquad a\longmapsto\sigma_a,\quad \sigma_a(\zeta_n)=\zeta_n^a
\tag{3}
\] is an isomorphism.

\textbf{Proof.} Let \(f\) be the monic minimal polynomial of
\(\zeta_n\). It divides the monic integral polynomial \(\Phi_n\); the
monic factor lemma makes its coefficients integers. Suppose \(p\nmid n\)
is prime and \(\zeta_n^p\) is not a root of \(f\). Let \(g\) be its
monic minimal polynomial, a distinct irreducible factor of \(\Phi_n\).
Since \(g(\zeta_n^p)=0\), we have \(f(X)\mid g(X^p)\) in
\(\mathbf Z[X]\).

Reduce modulo \(p\). The identity \(g(X^p)=g(X)^p\) there shows that any
irreducible factor of \(\overline f\) also divides \(\overline g\).
Since \(fg\mid X^n-1\), the latter polynomial would have a repeated
factor modulo \(p\). Its derivative \(nX^{n-1}\) is relatively prime to
it, which is a contradiction.

The same argument works for any primitive root that is already a root of
\(f\). Every positive integer \(a\) coprime to \(n\) is a product of
primes not dividing \(n\). Repeated application therefore puts every
\(\zeta_n^a\) among the roots of \(f\). Thus \(f=\Phi_n\). All these
roots lie in \(K_n\); sending \(\zeta_n\) to any one of them defines an
automorphism. Composition multiplies exponents, proving (3). The cases
\(n=1,2\) are linear polynomials and trivial groups. \(\square\)

The isomorphism (3) follows from the action on roots of unity.
Throughout this lesson Frobenius means the \textbf{arithmetic}
automorphism acting on residue fields by \(x\mapsto x^p\).

\subsection{References}\label{references}

The cyclotomic irreducibility argument is compared with J. S. Milne's
freely available
\href{https://www.jmilne.org/math/CourseNotes/FT.pdf}{Fields and Galois
Theory}, section on cyclotomic extensions, and
\href{https://www.jmilne.org/math/CourseNotes/ANT.pdf}{Algebraic Number
Theory}. The monic factor argument and complete Theorem 12.1 proof are
written above.

\end{document}
